\documentclass[11pt]{article}
\usepackage[hmargin=0.6in,vmargin=0.5in]{geometry}
\usepackage{amsmath}
\usepackage{amssymb}
\usepackage{amsthm}
\usepackage{microtype}
\usepackage[hidelinks]{hyperref}
\usepackage{enumitem}
\usepackage{longtable}
\usepackage{array}
\setlength{\parskip}{0pt}
\setlist{nosep,leftmargin=1.5em}

\newcommand{\R}{\mathbb{R}}
\newcommand{\N}{\mathbb{N}}
\newcommand{\supp}{\operatorname{supp}}
\newcommand{\Diff}{\operatorname{Diff}}
\newcommand{\Mo}{\mathring{M}}
\newcommand{\Mb}{\breve{M}}
\newcommand{\MB}{\boldsymbol M}
\newcommand{\xb}{\boldsymbol x}
\DeclareMathOperator{\ind}{ind}
\DeclareMathOperator{\dist}{dist}
\DeclareMathOperator{\interior}{int}
\newcommand{\Mstar}{(M$^*$)}

\title{\large Derivations accompanying ``A note on the topology-coincidence statements in\\ \emph{Topology of the space of conormal distributions}''}
\author{\normalsize Kunal Tyagi\thanks{\texttt{hello.jay.tyagi@gmail.com}. \textbf{Use of AI.} The mathematics in this note (the derivations, the computations, and the text) was produced by Claude (Anthropic) working under the author's direction; the author set the objectives, adjudicated the technical and editorial decisions, and is responsible for the content. Claude is a tool and is not an author.}}
\date{\vspace{-1ex}\normalsize September 6, 2026}

\begin{document}
\maketitle
\vspace{-2em}

\noindent This document accompanies the note \emph{A note on the topology-coincidence statements in ``Topology of the space of conormal distributions''} (below: the note). It contains full derivations of the witnesses (a)--(f) of the note's \S3, of the estimate \eqref{interp} of its \S4 (the note's (1)), and of the assembly of its \S4, with page anchors into the published text. The published paper \cite{ALKL24} is cited as [P] with its own printed page numbers (``Page $k$ of 68''); the memoir \cite{ALKL24m} is cited as [M] by section number, with the page numbers of its arXiv version (the same in v1 and v2). Statement numbers are those of the published article, which agree with arXiv v2 and v3 of [P].

\section*{0. Page anchors of every published statement quoted or cited in the note}
\noindent All entries were read in the published text of [P] (for the last three rows, in the arXiv text of [M]).

{\footnotesize
\begin{longtable}{@{}p{0.19\textwidth}p{0.07\textwidth}p{0.71\textwidth}@{}}
\hline
\textbf{Statement} & \textbf{Page} & \textbf{Verbatim (as extracted)}\\
\hline
\endhead
Acyclicity criterion & p.~4 & ``the condition of being acyclic can be described as follows [39, Theorem 6.1]: for all $k$, there is some $k'\ge k$ such that, for all $k''\ge k'$, the topologies of $X_{k'}$ and $X_{k''}$ coincide on some 0-neighborhood of $X_k$.''\\
regular / boundedly retractive; Prop.~6.4, Cor.~6.5 of [39] & pp.~4--5 & ``It is said that $(X_k)$ is regular if any bounded $B\subset X$ is contained and bounded in some step $X_k$. If moreover the topologies of $X$ and $X_k$ coincide on $B$, then $(X_k)$ is said to be boundedly retractive.'' \dots\ ``$X$ is acyclic if and only if it is boundedly/compactly/sequentially retractive [39, Proposition 6.4]. As a consequence, acyclic LF-spaces are complete and regular [39, Corollary 6.5].''\\
compact spectra & p.~5 & ``Assume the steps $X_k$ are LCHSs. It is said that $(X_k)$ is compact if the inclusion maps are compact operators. In this case, $(X_k)$ is clearly acyclic, and so $X$ is Hausdorff. Moreover $X$ is a complete bornological DF Montel space [18, Theorem 6$'$].''\\
(3.1) & p.~15 & $\|a\|_{K,\alpha,\beta,m}:=\sup_{x\in K,\,\xi\in\R^l}|\partial^\alpha_x\partial^\beta_\xi a(x,\xi)|/(1+|\xi|)^{m-|\beta|}<\infty$; ``a symbol of order at most $m\in\R$ on $U\times\R^l$''; ``$S^m(U\times\R^l)$, becomes a Fr\'echet space with the semi-norms'' (3.1).\\
(3.4), (3.5) & p.~16 & (3.4) $\|a\|_{Q,C^k}=\sup_{(x,\xi)\in Q,\ |\alpha|+|\beta|\le k}|\partial^\alpha_x\partial^\beta_\xi a|$; (3.5) $\|a\|'_{K,\alpha,\beta,m}=\sup_{x\in K}\limsup_{|\xi|\to\infty}|\partial^\alpha_x\partial^\beta_\xi a(x,\xi)|/|\xi|^{m-|\beta|}$.\\
Prop.~3.2 & p.~16 & ``The semi-norms (3.4) and (3.5) together describe the topology of $S^m(U\times\R^l)$.'' Proof pp.~16--17.\\
the failing sentence & p.~17, lines 1--2 & ``For any $a\in\hat S'^m(U\times\R^l)$, and $K$, $\alpha$ and $\beta$ like in (3.5), since $\|\varphi(a)\|'_{K,\alpha,\beta,m}=\|a\|'_{K,\alpha,\beta,m}<\infty$, there are $C,R>0$ so that, if $x\in K$ and $|\xi|\ge R$, then \dots'' and later ``Hence $S'^m(U\times\R^l)$ is complete, and therefore it is a Fr\'echet space.''\\
Prop.~3.3 & p.~17 & ``For $m,m'\in\N_0$ \dots\ if $m<m'$, then $\|\cdot\|'_{K,\alpha,\beta,m'}=0$ on $S^m(U\times\R^l)$.''\\
Cor.~3.4 & p.~17 & ``For $m<m'$, the topologies of $S^{m'}(U\times\R^l)$ and $C^\infty(U\times\R^l)$ coincide on $S^m(U\times\R^l)$. Therefore the topologies of $S^\infty(U\times\R^l)$ and $C^\infty(U\times\R^l)$ coincide on $S^m(U\times\R^l)$.'' Proof: ``The first assertion is a consequence of Propositions 3.2 and 3.3.''\\
Cor.~3.5 & p.~17 & density; ``Proof The first assertion is given by Corollary 3.4 and the density of $C_c^\infty(U\times\R^l)$ in $C^\infty(U\times\R^l)$.''\\
Cor.~3.6 & p.~18 & ``$S^\infty(U\times\R^l)$ is an acyclic Montel space, and therefore complete, boundedly retractive and reflexive. Proof Corollary 3.4 gives the property of being acyclic, and therefore complete and boundedly retractive (Sect. 2.1).'' Montel part: ``Since $S^\infty(U\times\R^l)$ is boundedly retractive, $B$ is contained and bounded in some $S^m(U\times\R^l)$ \dots\ By Corollary 3.4, it follows that $B$ is a complete bounded subspace of $C^\infty(U\times\R^l)$''.\\
Remark~3.7 & p.~18 & ``Another proof of Corollary 3.5 could be given like in Proposition 6.10.''\\
Remark~3.8 & p.~18 & uses ``since $S^\infty(U\times\R^l)$ is sequentially retractive (Corollary 3.6)''.\\
bundle extension & p.~18 & ``We can similarly define the norms (3.4) and (3.5) on $S^m(E)$, and Propositions 3.2 and 3.3 and Corollaries 3.4 to 3.6 can be directly extended to this setting.''\\
adapted chart & p.~19 & ``$(U,x)$ \dots\ adapted to $L$; i.e., for open subsets $U'\subset\R^{n'}$ and $U''\subset\R^{n''}$, $x=(x',x'')\colon U\to U'\times U''$ \dots\ $L_0:=L\cap U=\{x'=0\}$. If $L$ is of codimension one, then we will use the notation $(x,y)$ instead of $(x',x'')$.''\\
(4.7)--(4.8), Prop.~4.3 & p.~21 & $a(x'',\xi)=\int_{\R^{n'}}e^{-i\langle x',\xi\rangle}u(x',x'')\,dx'$; $u(x)=(2\pi)^{-n'}\int e^{i\langle x',\xi\rangle}a(x'',\xi)\,d\xi$.\\
\S4.3.2, (4.9), (4.10), (4.12) & p.~22 & ``Take a finite cover of $L$ by relatively compact charts $(U_j,x_j)$ of $M$ adapted to $L$, and write $L_j=L\cap U_j$. Let $\{h,f_j\}$ be a $C^\infty$ partition of unity of $M$ subordinated to the open covering $\{M\setminus L,U_j\}$''; (4.9) $\bar m=m+n/4-n'/2$; ``$I^m(M,L)$, which becomes a Fr\'echet space''; ``the following map is required to be a TVS-embedding: (4.10) $I^m(M,L)\to C^\infty(M\setminus L)\oplus\bigoplus_jS^{\bar m}(N^*L_j;\Omega N^*L_j)$, $u\mapsto(hu,(a_j))$''; (4.12) $I^{(-m-n/4+\epsilon)}\subset I^m\subset I^{(-m-n/4-\epsilon)}$.\\
Cor.~4.5 & p.~23 & ``For $m<m',m''$, the topologies of $I^{m'}(M,L)$ and $I^{m''}(M,L)$ coincide on $I^m(M,L)$. Proof Use Corollary 3.4 and the TVS-embeddings (4.10).''\\
Cor.~4.6 & p.~23 & ``$C^\infty(M)$ is contained in the stated spaces by (4.5).''\\
Cor.~4.7 & p.~23 & ``$I(M,L)$ is an acyclic Montel space \dots\ Proof Like in Corollary 3.6, by Corollaries 4.2 and 4.5, it is enough to prove that $I(M,L)$ is semi-Montel.'' \dots\ ``Then $I(M,L)$ is semi-Montel because $C^\infty(M\setminus L)$ and $S^\infty(N^*L_j;\Omega N^*L_j)$ are Montel spaces (Corollary 3.6)''.\\
\S4.3.3 & p.~23 & ``Corollary 4.7 has extensions for $\bigcup_mI^m(M,L)$ and $I_{\cdot/c}(M,L)$, except acyclicity in the case of $I(M,L)$.''\\
$\Diff_b$ & p.~35 & ``$\mathfrak X_b(M)\subset\mathfrak X(M)$ of vector fields tangent to $\partial M$, called $b$-vector fields \dots\ $\Diff_b(M)\subset\Diff(M)$ of $b$-differential operators.''\\
(6.33) & p.~36 & ``$\dot A(M)|_{\Mo},A(M)\subset C^\infty(\Mo)$'' continuous inclusions.\\
$x^mL^\infty(M)$ & p.~37 & ``there is a continuous inclusion $x^mL^\infty(M)\subset C^{-\infty}(M)$.''\\
$A^m(M)$, (6.38), (6.40), (6.41) & p.~38 & $A^m(M)=\{u\in C^{-\infty}(M)\mid\Diff_b(M)u\subset x^mL^\infty(M)\}$, ``projective topology given by the maps $P\colon A^m(M)\to x^mL^\infty(M)$ ($P\in\Diff_b(M)$)''; (6.38) $A^m\subset A^{m'}$ ($m'<m$); (6.40) $A(M)=\bigcup_mA^m(M)$; ``Let $\{P_j\mid j\in\N_0\}$ be a countable $C^\infty(M)$-spanning set of $\Diff_b(M)$. The topology of $A^m(M)$ can be described by the semi-norms $\|\cdot\|_{k,m}$ ($k\in\N_0$)''; (6.41) $\|u\|_{k,m}=\sup_{\Mo}x^{-m}|P_ku|$.\\
(6.42), (6.43), Prop.~6.12, Cor.~6.14, Cor.~6.16 & p.~39 & (6.42) $\|u\|_{K,k,m}=\sup_K|P_ku|$; (6.43) $\|u\|'_{k,m}=\lim_{\epsilon\downarrow0}\sup_{\{0<x<\epsilon\}}|x^{-m}P_ku|$; ``The proofs of the following results are similar to the proofs of Propositions 3.2 and 3.3, Corollaries 3.4 to 3.6, using (6.33).'' Prop.~6.12 ``The semi-norms (6.42) and (6.43) together describe the topology of $A^m(M)$.'' Cor.~6.14 ``If $m'<m$, then the topologies of $A^{m'}(M)$ and $C^\infty(\Mo)$ coincide on $A^m(M)$. Therefore the topologies of $A(M)$ and $C^\infty(\Mo)$ coincide on $A^m(M)$.'' Cor.~6.16 ``$A(M)$ is an acyclic Montel space\dots''.\\
(6.47), Cor.~6.21, Cor.~6.22, $\mathcal K^m(M)$ & p.~40 & (6.47) $\dot A^m(M)=I^m_M(\Mb,\partial M)\subset I^m(\Mb,\partial M)$, ``closed subspaces''; ``The following is a consequence of Corollary 4.5 applied to $(\Mb,\partial M)$. Corollary 6.21 For $m<m',m''$, the topologies of $\dot A^{m'}(M)$ and $\dot A^{m''}(M)$ coincide on $\dot A^m(M)$.'' Cor.~6.22 ``follows like Corollary 3.6, applying Corollary 6.21''. Sect.~6.13: ``$\mathcal K^m(M)=\dot A^m_{\partial M}(M)$''.\\
$\mathcal K^m(M)$, Cor.~6.27, Cor.~6.28, (6.49) & p.~41 & $\mathcal K^{(s)},\mathcal K^m,\mathcal K$ ``closed subspaces of $\dot A^{(s)}(M)$, $\dot A^m(M)$ and $\dot A(M)$ \dots\ null spaces of the corresponding restrictions of the map (6.36)''; ``the following analogs of Corollaries 6.21 and 6.22 hold true with formally the same proofs, using Corollaries 6.21, 6.22 and 6.26. Corollary 6.27 For $m<m',m''$, the topologies of $\mathcal K^{m'}(M)$ and $\mathcal K^{m''}(M)$ coincide on $\mathcal K^m(M)$.'' ``By Corollary 6.20, (6.49) $\mathcal K(M)\equiv I_{\partial M}(\Mb,\partial M)$''.\\
Prop.~6.29 & p.~42 & ``For all $m\in\R$, there is a continuous linear partial extension map $E_m\colon A^m(M)\to\dot A^{(s)}(M)$, where $s=0$ if $m\ge0$, and $m>s\in\mathbb Z^-$ if $m<0$.'' Definition above it: ``a map $E\colon X\to Y$ is called a partial extension map if $R(Y)\subset X$ and $RE=1$ on $X$.''\\
Cor.~6.38, Remark~6.41 & pp.~46--47 & ``Corollary 6.38 $A^m(M)\equiv x^mH_b^\infty(M)\equiv x^{m+1/2}H^\infty(\Mo)$ ($m\in\R$)''; ``Remark 6.41 Corollary 6.39 and the first identities of Corollaries 6.38 and 6.40 are independent of $g$. So they hold true without the assumptions (A) and (B).''\\
Prop.~6.45, Claim~6.46 and its proof & p.~48 & Claim 6.46 ``For all bounded subset $A\subset A(M)$, there is some bounded subset $B\subset\dot A(M)$ such that, for all 0-neighborhood $U\subset\dot A(M)$, there is a 0-neighborhood $V\subset A(M)$ so that $A\cap V\subset R(B\cap U)$.'' Proof: ``Since $A(M)$ is boundedly retractive (Corollary 6.16), $A$ is contained and bounded in some step $A^m(M)$. For any $m'>m$, let $E_{m'}\colon A^m(M)\to\dot A^{(s)}(M)$ be the partial extension map given by Proposition 6.29. Then $B:=E_{m'}(A)$ is bounded \dots\ given any 0-neighborhood $U\subset\dot A(M)$, there is some 0-neighborhood $W\subset A^{m'}(M)$ so that $E_{m'}(W)\subset U\cap\dot A^{(s)}(M)$. By Corollary 6.14, [p.~49:] there is some 0-neighborhood $V\subset A(M)$ such that $V\cap A^m(M)=W\cap A^m(M)$.''\\
cutting along $L$ & p.~50 & ``Let $\MB$ be the smooth manifold with boundary defined by `cutting' $M$ along $L$''.\\
(7.26), (7.27), (7.28), Cor.~7.13 & p.~56 & (7.26) $\pi_*\colon A(\MB)\to J(M,L)$ TVS-isomorphism; ``We also get spaces $J^{(s)}(M,L)$ and $J^m(M,L)$ \dots\ corresponding to $A^{(s)}(\MB)$ and $A^m(\MB)$ via (7.26). Extend $|x|$ to a function $\xb$ on $M$ that is positive and smooth on $M\setminus L$. Its lift $\pi^*\xb$ is a boundary defining function of $\MB$, also denoted by $\xb$''; (7.27) $J^m(M,L)=\{u\in C^{-\infty}(M,L)\mid\Diff(M,L)u\subset\xb^mL^\infty(M)\}$; (7.28) $C^\infty(M)\subset J^{(\infty)}(M,L)$; Cor.~7.13 ``If $m'<m$, then the topologies of $J^{m'}(M,L)$ and $C^\infty(M\setminus L)$ coincide on $J^m(M,L)$. Therefore the topologies of $J(M,L)$ and $C^\infty(M\setminus L)$ coincide on $J^m(M,L)$.''\\
Cor.~7.15 & p.~57 & ``$J(M,L)$ is an acyclic Montel space\dots''\\
$K^m(M,L)$, Cor.~7.22, Cor.~7.23 & p.~58 & ``$K^m(M,L)=I^m_L(M,L)$ \dots\ closed subspaces \dots\ null spaces of the corresponding restrictions of the map (7.31)''; ``we get the following consequences of Propositions 6.24 and 6.25 and Corollaries 6.26 to 6.28 \dots\ Corollary 7.22 For $m<m',m''$, the topologies of $K^{m'}(M,L)$ and $K^{m''}(M,L)$ coincide on $K^m(M,L)$.''\\
Prop.~7.26 & p.~59 & ``The map (7.36) is a TVS-isomorphism'' ($K(M,L)$ as the locally convex direct sum $\bigoplus_mC^\infty(L;\Omega^{-1}NL)$).\\
Cor.~7.31 & p.~60 & ``For all $m\in\R$, there is a continuous linear partial extension map $E_m\colon J^m(M,L)\to I^{(s)}(M,L)$''.\\
Prop.~8.8 & p.~64 & ``The following analog of Proposition 6.45 holds true with formally the same proof, using Proposition 7.29 and Corollaries 7.13, 7.15 and 7.31. Proposition 8.8 The dual-conormal sequence of $M$ at $L$ is exact''.\\
refs [18], [39] & pp.~67--68 & ``18. Komatsu, H.: Projective and injective limits of weakly compact sequences of locally convex spaces. J. Math. Soc. Japan 19, 366--383 (1967)''; ``39. Wengenroth, J.: Derived Functors in Functional Analysis. Lecture Notes in Mathematics, vol. 1810. Springer, Berlin (2003)''.\\
{[M]} \S2.1.8 & arXiv p.~15 & ``The following properties hold [\'ALKL23, Corollaries 3.4--3.6 and Remark 3.8]: The topologies of $S^\infty(U\times\R^l)$ and $C^\infty(U\times\R^l)$ coincide on $S^m(U\times\R^l)$, however the second inclusion of (2.1.27) is not a TVS-embedding; $C_c^\infty(U\times\R^l)$ is dense in $S^\infty(U\times\R^l)$; and $S^\infty(U\times\R^l)$ is an acyclic Montel space, and therefore complete, boundedly/compactly/sequentially retractive and reflexive.'' (``For any open $U\subset\R^n$''.)\\
{[M]} \S2.5.10, \S2.6.7 & arXiv pp.~38, 53 & ``the topologies of $A(M)$ and $C^\infty(\Mo)$ coincide on every $A^m(M)$''; ``the topologies of $J(M,L)$ and $C^\infty(M\setminus L)$ coincide on every $J^m(M,L)$''.\\
{[M]} consumers & arXiv pp.~119, 122 & \S5.2.1 ``because $I(\mathcal F)$ is compactly retractive (Section 2.2.2)''; \S5.5.3 ``Using that $J(\mathcal F)$ is compactly retractive (Section 2.6.7)''; \S5.5.4 ``Since $I(\mathcal F)$ is compactly retractive''. $I(\mathcal F):=I(M,M^0;\Lambda\mathcal F)$ (arXiv p.~117), $M$ closed.\\
\hline
\end{longtable}
}

\noindent\textbf{Numbering across versions.} arXiv v1 of [P] has Prop.~6.10 / Cor.~6.12 / Cor.~6.19 / Cor.~6.24 / Prop.~6.26 / Claim~6.43 / Cor.~7.11 / Cor.~7.20 / Cor.~7.28 for the statements printed as Prop.~6.12 / Cor.~6.14 / Cor.~6.21 / Cor.~6.27 / Prop.~6.29 / Claim~6.46 / Cor.~7.13 / Cor.~7.22 / Cor.~7.31 in the published text; the numbers of \S\S3--4 are identical in all versions. The note and this document use the published numbers only.

\section*{1. Conventions}
$U\subset\R^n$ open ($n\ge0$), $l\ge1$. $S^m=S^m(U\times\R^l)$ with the seminorms (3.1). Every (3.4)- and (3.5)-seminorm is continuous on $S^m$: $\|a\|_{Q,C^k}\le\max_{|\alpha|+|\beta|\le k}\sup_Q(1+|\xi|)^{m-|\beta|}\,\|a\|_{K,\alpha,\beta,m}$ for $Q\subset K\times\bar B(0,\rho)$, and $\|a\|'_{K,\alpha,\beta,m}\le\|a\|_{K,\alpha,\beta,m}$ because $|\xi|^{|\beta|-m}/(1+|\xi|)^{|\beta|-m}\to1$. So the topology $S'^m$ of (3.4)+(3.5) is coarser than or equal to that of $S^m$, and the $C^\infty$-topology (given by (3.4)) is coarser than or equal to $S'^m$. $\theta\in C_c^\infty(\R^l)$ is a fixed bump with $\theta(0)=1$ and $\supp\theta\subset B(0,r)$. All witnesses are independent of $x$ (or have a fixed compactly supported factor in the base variable).

\section*{2. Witness (a): $g_N$. Prop.~3.2 and the first assertion of Cor.~3.4 are false}
$g_N(x,\xi):=\theta(\xi-Ne_1)$, $N\in\N$.
\begin{enumerate}
\item $g_N\in C^\infty(U\times\R^l)$, independent of $x$, compactly supported in $\xi$, so $\|g_N\|_{K,\alpha,\beta,m}<\infty$ for all $\alpha,\beta,m$ ($\alpha\ne0$ gives $0$; $\alpha=0$: a bounded function times a bounded weight on a compact set). Hence $g_N\in S^{-\infty}\subset S^m$ for every $m\in\R$.
\item (3.4): if $Q\subset K\times\bar B(0,\rho)$ then $\supp g_N\subset U\times\{|\xi|\ge N-r\}$, so $\|g_N\|_{Q,C^k}=0$ for $N>\rho+r$. Hence $g_N\to0$ in $C^\infty(U\times\R^l)$.
\item (3.5): for each $x$, $\partial^\alpha_x\partial^\beta_\xi g_N(x,\xi)=0$ for $|\xi|>N+r$, so the limsup is $0$: every (3.5)-seminorm of $g_N$ is $0$.
\item (3.1) with $\alpha=0$: $\|g_N\|_{K,0,\beta,m}=\sup_\xi|\partial^\beta\theta(\xi-Ne_1)|(1+|\xi|)^{|\beta|-m}=\sup_{|\eta|\le r}|\partial^\beta\theta(\eta)|(1+|\eta+Ne_1|)^{|\beta|-m}$. For $|\eta|\le r$, $N-r\le|\eta+Ne_1|\le N+r$. If $|\beta|-m>0$, then $\|g_N\|_{K,0,\beta,m}\ge(1+N-r)^{|\beta|-m}\sup|\partial^\beta\theta|\to\infty$ provided $\partial^\beta\theta\not\equiv0$.
\item Existence of $\beta$ with $|\beta|>m$ and $\partial^\beta\theta\not\equiv0$: if all derivatives of some order $k$ vanished identically, $\theta$ would be a polynomial of degree $<k$, and a compactly supported polynomial is $0$, contradicting $\theta(0)=1$. So for every $k$ there is $\beta$ with $|\beta|=k$ and $\partial^\beta\theta\not\equiv0$.
\end{enumerate}
Conclusions. (i) $g_N\to0$ in $S'^m$ but $(g_N)$ is unbounded in $S^m$, so the identity $S^m\to S'^m$ is continuous and not a homeomorphism: \textbf{Prop.~3.2 is false} for every $m\in\R$, every $U$, every $l\ge1$. (ii) For $m<m'$ choose $|\beta|>m'$: $g_N\in S^m$, $g_N\to0$ in $C^\infty$, $g_N\not\to0$ in $S^{m'}$ (unbounded). \textbf{The first assertion of Cor.~3.4 is false} for every $m<m'$. (In fact for all $m,m'$.)

\section*{3. Witness (b): a Cauchy sequence of $S'^m$ with no limit; the failing sentence}
Fix $R>2r$, so the translates $\theta(\xi-iRe_1)$, $i\in\N$, have pairwise disjoint supports. $c_i:=(1+iR)^{m+1}$, $b_j:=\sum_{i=1}^jc_i\theta(\xi-iRe_1)$ (independent of $x$). Each $b_j\in S^{-\infty}\subset S^m$.
\begin{enumerate}
\item For $i<j$, $b_j-b_i=\sum_{k=i+1}^jc_k\theta(\xi-kRe_1)$ is supported in $\{|\xi|\ge(i+1)R-r\}$ and has compact $\xi$-support. So all its (3.5)-seminorms vanish, and for $Q\subset K\times\bar B(0,\rho)$ its (3.4)-seminorms vanish as soon as $(i+1)R-r>\rho$. Hence for every continuous seminorm $p$ of $S'^m$ there is $i_0$ with $p(b_j-b_i)=0$ for all $j>i\ge i_0$: \textbf{$(b_j)$ is Cauchy in $S'^m$.}
\item $b:=\sum_{i\ge1}c_i\theta(\xi-iRe_1)$ is a locally finite sum, hence smooth, and $b_j\to b$ in $C^\infty(U\times\R^l)$ (on a compact $\xi$-set, $b_j=b$ for $j$ large).
\item $b\notin S^m$: at $\xi=iRe_1$ only the $i$-th term is nonzero, so $|b(x,iRe_1)|=c_i$, and $|b(x,iRe_1)|(1+|iRe_1|)^{-m}=(1+iR)^{m+1}(1+iR)^{-m}=1+iR\to\infty$; thus $\|b\|_{K,0,0,m}=\infty$ (any real $m$).
\item If $b_j\to\tilde b$ in $S'^m$ with $\tilde b\in S^m$, then $b_j\to\tilde b$ in $C^\infty$ ($S'^m\to C^\infty$ continuous), so $\tilde b=b$ by Hausdorffness of $C^\infty$, contradicting 3. \textbf{So $(b_j)$ has no limit in $S'^m$; $S'^m$ is not complete}, contrary to ``Hence $S'^m(U\times\R^l)$ is complete'' (p.~17).
\item The failing sentence. Let $a\in\hat S'^m$ be the class of $(b_j)$ in the completion. The continuous extension of $\|\cdot\|'_{K,0,0,m}$ satisfies $\|a\|'_{K,0,0,m}=\lim_j\|b_j\|'_{K,0,0,m}=0$. The extended map $\varphi\colon\hat S'^m\to C^\infty$ sends $a$ to $\lim b_j=b$, and $\|\varphi(a)\|'_{K,0,0,m}=\|b\|'_{K,0,0,m}=\sup_x\limsup_{|\xi|\to\infty}|b(x,\xi)|/|\xi|^m\ge\limsup_i(1+iR)^{m+1}/(iR)^m=\infty$. So the printed ``$\|\varphi(a)\|'_{K,\alpha,\beta,m}=\|a\|'_{K,\alpha,\beta,m}<\infty$'' (p.~17, line 2) is false: the (3.5)-seminorms of the completion are not computed by $\varphi$.
\end{enumerate}
(A second gap in the same sentence: finiteness of a sup over $x\in K$ of limsups gives no threshold $R$ uniform in $x\in K$. Not used in the note.)

\section*{4. Witness (c): the second assertion of Cor.~3.4 is false}
$c_j:=e^j\theta(\xi-je_1)$, $j\in\N$. As in \S2, $c_j\in S^{-\infty}\subset S^m$ for all $m$, and $c_j\to0$ in $C^\infty$ (supports leave every compact set).

Fix $x_0\in U$, $K:=\{x_0\}$. For $k\in\N$: $\varepsilon_k:=\tfrac12\inf_{j\ge1}e^j(1+j)^{-k}$. The sequence $e^j(1+j)^{-k}$ is positive and tends to $\infty$, so $\varepsilon_k>0$. $W_k:=\{a\in S^k:\|a\|_{K,0,0,k}<\varepsilon_k\}$, an open 0-neighborhood of $S^k$. $W:=$ absolutely convex hull of $\bigcup_{k\ge1}W_k$ (finite sums $\sum\lambda_iw_i$ with $\sum|\lambda_i|\le1$, $w_i\in\bigcup W_k$).
\begin{enumerate}
\item $W$ is a 0-neighborhood of $S^\infty=\ind_kS^k$: it is absolutely convex and $W\cap S^k\supseteq W_k$ is a 0-neighborhood of $S^k$ for every $k\in\N$ (definition of the locally convex inductive limit topology: the absolutely convex sets meeting every step in a 0-neighborhood form a 0-basis). Since $S^\infty=\bigcup_{m\in\R}S^m=\bigcup_{k\in\N}S^k$ with continuous inclusions $S^m\subset S^k$ ($k\ge m$), the inductive limit over $\R$ and over $\N$ carry the same topology (the paper's own real-index license, p.~5).
\item If $c_j\in W$, write $c_j=\sum_i\lambda_iw_i$ as above. For $w\in W_k$, $|w(x_0,\xi)|\le\|w\|_{K,0,0,k}(1+|\xi|)^k<\varepsilon_k(1+|\xi|)^k$. At $\xi=je_1$: $|w_i(x_0,je_1)|<\varepsilon_{k_i}(1+j)^{k_i}\le\tfrac12e^j$ (by definition of $\varepsilon_{k_i}$, $\varepsilon_{k_i}(1+j)^{k_i}\le\tfrac12e^j(1+j)^{-k_i}(1+j)^{k_i}$). Hence $e^j=|c_j(x_0,je_1)|\le\sum|\lambda_i|\,|w_i(x_0,je_1)|<\tfrac12e^j$, contradiction.
\item So $c_j\notin W$ for every $j$: $c_j\not\to0$ in $S^\infty$, while $c_j\to0$ in $C^\infty$ and $c_j\in S^m$ for every $m$. \textbf{The topologies of $S^\infty$ and $C^\infty$ differ on every $S^m$.} No regularity of $S^\infty$ is used.
\end{enumerate}

\section*{5. Witness (d): Cor.~4.5 is false for every compact $(M,L)$, $\operatorname{codim}\ge1$; Cor.~6.21}
Setting: $M$ compact, $L$ closed regular submanifold of codimension $n'\ge1$, $n=\dim M$, $n''=n-n'$. Chart $(U_1,x=(x',x''))$ adapted to $L$ (p.~19: $L\cap U_1=\{x'=0\}$, $x'\in\R^{n'}$), chosen around $p_0\in L$ with $x(p_0)=0$ and $x(U_1)\supset\{|x'|<1\}\times U''$. $m<m'<m''$, $\bar m=m+n/4-n'/2$ etc.\ (4.9); $\bar m'<\bar m''$.

Partition of unity: given any $\{h,f_j\}$ subordinate to $\{M\setminus L,U_j\}$ and $\rho\in C_c^\infty(U_1)$, $0\le\rho\le1$, with $\rho=1$ on a neighborhood $N$ of the compact set $x^{-1}(\{0\}\times\supp g)\subset L$ (with $g$ as below), the family $\tilde h:=(1-\rho)h$, $\tilde f_1:=\rho+(1-\rho)f_1$, $\tilde f_j:=(1-\rho)f_j$ ($j\ge2$) is smooth, nonnegative, sums to $\rho+(1-\rho)=1$, has $\supp\tilde h\subset\supp h\subset M\setminus L$, $\supp\tilde f_j\subset\supp f_j\subset U_j$ ($j\ge2$), $\supp\tilde f_1\subset\supp\rho\cup\supp f_1\subset U_1$: a partition of unity subordinate to the same cover with $\tilde f_1=1$, $\tilde h=\tilde f_j=0$ on $N$. The paper allows any such partition (p.~22). Independence of the topology of $I^m$ from this choice: any two Fr\'echet topologies on $I^m(M,L)$ both finer than the (Hausdorff) topology induced from $C^{-\infty}(M)$ coincide by the closed graph theorem ($I^m\subset I^{(s)}\subset H^s(M)\subset C^{-\infty}(M)$ continuously by (4.12) and the definitions). So the refutation is choice-independent.

Data: $g\in C_c^\infty(U'')$, $g\not\equiv0$; $\psi\in C_c^\infty(\R^{n'})$ with $\supp\psi\subset\{|x'|<1\}$, $\psi\not\equiv0$; $N_0\in\N_0$ with $2N_0\ge\bar m''$; $\psi_*:=\Delta^{N_0}\psi$ (Laplacian in $x'$). Then $\hat\psi_*(\eta)=(-|\eta|^2)^{N_0}\hat\psi(\eta)$ (Fourier convention (4.8): $a(\xi)=\int e^{-i\langle x',\xi\rangle}u\,dx'$; $\partial_{x'_i}\mapsto i\xi_i$, $\Delta\mapsto-|\xi|^2$), Schwartz, $\not\equiv0$, and for $|\beta|\le2N_0$ and $|\eta|\le1$: $|\partial^\beta\hat\psi_*(\eta)|\le C_\beta|\eta|^{2N_0-|\beta|}$ (Leibniz: derivatives of the homogeneous polynomial $|\eta|^{2N_0}$ of order $|\gamma|\le|\beta|$ are $O(|\eta|^{2N_0-|\gamma|})=O(|\eta|^{2N_0-|\beta|})$ on $|\eta|\le1$).

$u_j(x',x''):=c_j\,g(x'')\,\psi_*(jx')$, $c_j:=j^{n'+\bar m'}$, $j\ge1$. $\supp u_j\subset x^{-1}(\{|x'|\le1/j\}\times\supp g)\subset N$ for $j\ge j_0$. $u_j\in C_c^\infty(U_1)\subset C^\infty(M)\subset I^{m_1}(M,L)$ for every $m_1$ (p.~23 ``$C^\infty(M)$ is contained in the stated spaces by (4.5)''; equivalently, its (4.10)-image below lies in every $S^{\bar m_1}$).

(4.10)-image for $j\ge j_0$: $\tilde hu_j=0$, $\tilde f_ju_j=0$ ($j\ge2$), $\tilde f_1u_j=u_j$, whose symbol by (4.8) is
\[
a_j(x'',\xi)=\int e^{-i\langle x',\xi\rangle}c_j\,g(x'')\,\psi_*(jx')\,dx'=c_j\,j^{-n'}g(x'')\,\hat\psi_*(\xi/j).
\]
Seminorms (3.1) over a compact $K\subset U''$ ($K\supseteq\supp g$): with $\xi=j\eta$,
\[
|\partial^\alpha_{x''}\partial^\beta_\xi a_j|(1+|\xi|)^{|\beta|-\bar m_1}=c_jj^{-n'}|\partial^\alpha g(x'')|\,j^{-|\beta|}|\partial^\beta\hat\psi_*(\eta)|\,(1+j|\eta|)^{|\beta|-\bar m_1}=c_jj^{-n'-\bar m_1}|\partial^\alpha g(x'')|\,|\partial^\beta\hat\psi_*(\eta)|\,(1/j+|\eta|)^{|\beta|-\bar m_1},
\]
using $j^{-|\beta|}(1+j|\eta|)^{|\beta|-\bar m_1}=j^{-|\beta|}j^{|\beta|-\bar m_1}(1/j+|\eta|)^{|\beta|-\bar m_1}$. Hence
\[
\|a_j\|_{K,\alpha,\beta,\bar m_1}=c_jj^{-n'-\bar m_1}\|\partial^\alpha g\|_{L^\infty(K)}\,\Phi_{\beta,\bar m_1}(j),\qquad\Phi_{\beta,\bar m_1}(j):=\sup_\eta|\partial^\beta\hat\psi_*(\eta)|\,(1/j+|\eta|)^{|\beta|-\bar m_1}.
\]
Upper bound ($\bar m_1\le2N_0$): if $|\beta|\ge\bar m_1$, $(1/j+|\eta|)^{|\beta|-\bar m_1}\le(1+|\eta|)^{|\beta|-\bar m_1}$ and $\hat\psi_*$ Schwartz give $\Phi\le C$. If $|\beta|<\bar m_1\le2N_0$: $(1/j+|\eta|)^{|\beta|-\bar m_1}\le|\eta|^{|\beta|-\bar m_1}$; for $|\eta|\le1$, $|\partial^\beta\hat\psi_*(\eta)||\eta|^{|\beta|-\bar m_1}\le C_\beta|\eta|^{2N_0-\bar m_1}\le C_\beta$; for $|\eta|\ge1$ the weight is $\le1$ and $\partial^\beta\hat\psi_*$ is bounded. So $\Phi_{\beta,\bar m_1}(j)\le C_{\beta,\bar m_1}$ for all $j\ge1$.

Lower bound: pick $\eta_0\ne0$ with $\hat\psi_*(\eta_0)\ne0$ (exists: $\hat\psi_*$ is entire and $\not\equiv0$, so its zero set has empty interior). $\Phi_{0,\bar m'}(j)\ge|\hat\psi_*(\eta_0)|(1/j+|\eta_0|)^{-\bar m'}\ge|\hat\psi_*(\eta_0)|\min\{(1+|\eta_0|)^{-\bar m'},|\eta_0|^{-\bar m'}\}=:c>0$ for all $j\ge1$.

With $c_j=j^{n'+\bar m'}$: $\|a_j\|_{K,\alpha,\beta,\bar m''}\le Cj^{\bar m'-\bar m''}\to0$ for every $\alpha,\beta,K$ ($\bar m'<\bar m''\le2N_0$), so \textbf{$u_j\to0$ in $I^{m''}(M,L)$} (its topology is the initial topology of (4.10)); and $\|a_j\|_{K,0,0,\bar m'}\ge c\|g\|_\infty>0$, so \textbf{$u_j\not\to0$ in $I^{m'}(M,L)$}. Since $u_j\in I^m(M,L)$, the $I^{m'}$- and $I^{m''}$-topologies differ on $I^m(M,L)$: \textbf{Cor.~4.5 is false} for every compact $(M,L)$ with $\operatorname{codim}L\ge1$ and every $m<m'<m''$ (the statement is symmetric in $m',m''$).

Cor.~6.21 (p.~40): $\dot A^m(M)=I^m_M(\Mb,\partial M)\subset I^m(\Mb,\partial M)$ closed subspaces (6.47), induced topologies. Run the above on $(\Mb,\partial M)$ ($n'=1$, chart $(x,y)$, $M=\{x\ge0\}$) with $\supp\psi\subset(\tfrac12,1)$; then $\psi_*=\psi^{(2N_0)}$ has support in $(\tfrac12,1)$, and $\supp u_j\subset\{1/(2j)\le x\le1/j\}\times\supp g\subset\Mo$, so $u_j\in C_c^\infty(\Mo)\subset\dot A^{m_1}(M)$ for all $m_1$, $u_j\to0$ in $\dot A^{m''}(M)$ and $u_j\not\to0$ in $\dot A^{m'}(M)$. \textbf{Cor.~6.21 is false.}

\section*{6. Witness (e): Prop.~6.12, Cor.~6.14, Cor.~7.13}
Setting (\S6 of [P]): $M$ compact with boundary, $x$ a boundary defining function, collar chart $(x,y)\in[0,x_0)\times V_1$, $V_1\subset\R^{n-1}$. $\Diff_b(M)$: generated by vector fields tangent to $\partial M$ (p.~35); on the collar these are $C^\infty$-combinations of $x\partial_x$ and $\partial_{y_i}$ (a vector field $a\partial_x+\sum b_i\partial_{y_i}$ is tangent to $\{x=0\}$ iff $a$ vanishes there iff $a=x\tilde a$). Topology of $A^m(M)$: the seminorms $\sup_{\Mo}x^{-m}|Pu|$, $P\in\Diff_b(M)$ (p.~38); in particular $P=1$ gives the continuous seminorm $\sup_{\Mo}x^{-m}|u|$. (6.33): $A(M)\subset C^\infty(\Mo)$ continuously, so pointwise evaluation on $\Mo$ is legitimate.

$\chi\in C_c^\infty((\tfrac12,2))$, $\chi(1)=1$; $g\in C_c^\infty(V_1)$, $g(y_0)=1$; $m'<m$.

$u_j:=j^{-m'}\chi(jx)g(y)$. $\supp u_j\subset\{1/(2j)<x<2/j\}\times\supp g$, compact in $\Mo$ for $2/j<x_0$.
\begin{enumerate}
\item $u_j\in C_c^\infty(\Mo)\subset A^{m_1}(M)$ for every $m_1$: for $P\in\Diff_b(M)$ (or any differential operator with smooth coefficients), $Pu_j$ is smooth with compact support in $\Mo$, where $x$ is bounded below, so $x^{-m_1}Pu_j$ is bounded.
\item $u_j\to0$ in $C^\infty(\Mo)$: on any compact $K\subset\Mo$, $u_j=0$ for $j$ large ($x\ge\min_Kx>2/j$).
\item (6.42): $\sup_K|P_ku_j|=0$ for $j$ large, by 2. (6.43): $u_j=0$ on $\{x<1/(2j)\}$, so $\lim_{\epsilon\downarrow0}\sup_{0<x<\epsilon}x^{-m}|P_ku_j|=0$ for every $j,m,k$.
\item (6.41) with $P=1$ and order $m'$: $\sup_{\Mo}x^{-m'}|u_j|\ge x^{-m'}|u_j|$ at $(1/j,y_0)$ $=j^{m'}\cdot j^{-m'}\chi(1)g(y_0)=1$.
\end{enumerate}
Conclusions. With $m'=m$: $u_j\to0$ in the (6.42)+(6.43) topology of $A^m(M)$ but $u_j\not\to0$ in $A^m(M)$: \textbf{Prop.~6.12 is false.} With $m'<m$: $u_j\in A^m(M)$, $u_j\to0$ in $C^\infty(\Mo)$, $u_j\not\to0$ in $A^{m'}(M)$: \textbf{the first assertion of Cor.~6.14 is false.}

Second assertion: $v_j:=e^j\chi(jx)g(y)\in A^{m_1}(M)$ for all $m_1$, $v_j\to0$ in $C^\infty(\Mo)$. $A(M)=\bigcup_{m\in\R}A^m(M)$ (6.40) $=\bigcup_{k\in\N}A^{-k}(M)$ by (6.38) ($A^m\subset A^{-k}$ for $-k<m$), same inductive topology. $\varepsilon_k:=\tfrac12\inf_{i\ge1}e^ii^{-k}>0$; $W_k:=\{u\in A^{-k}(M):\sup_{\Mo}x^k|u|<\varepsilon_k\}$ (a 0-neighborhood of $A^{-k}$: $P=1$, order $-k$); $W:=$ absolutely convex hull of $\bigcup_kW_k$, a 0-neighborhood of $A(M)$. If $v_j=\sum\lambda_iw_i$ with $\sum|\lambda_i|\le1$, $w_i\in W_{k_i}$: $|w_i(1/j,y_0)|\le\varepsilon_{k_i}(1/j)^{-k_i}=\varepsilon_{k_i}j^{k_i}\le\tfrac12e^j$, so $e^j=|v_j(1/j,y_0)|<\tfrac12e^j$, contradiction. So $v_j\notin W$ for all $j$: \textbf{$v_j\not\to0$ in $A(M)$}, the second assertion of Cor.~6.14 is false. (Consequently, no $V$ with $V\cap A^m=W\cap A^m$ as in the printed Claim~6.46 exists in general: the $A^{m'}$- and $A^{m''}$-topologies, $m''<m'<m$, differ on $A^m$, since $x^{-m''}|(x\partial_x)^a\partial_y^bu_j|\lesssim j^{m''-m'}\to0$ while $\sup x^{-m'}|u_j|\ge1$.)

Cor.~7.13 (p.~56): $M$ closed, $L$ a hypersurface, $J^m(M,L)=\{u\in C^{-\infty}(M,L):\Diff(M,L)u\subset\xb^mL^\infty(M)\}$ (7.27), topology ``like in Sects.~6.8 and 6.10'', i.e.\ the seminorms $\sup\xb^{-m}|Pu|$, $P\in\Diff(M,L)$ (again $P=1$ allowed); $\xb$ extends $|x|$ of an adapted chart $(x,y)$ with $L=\{x=0\}$ (p.~56), and is bounded below on compact subsets of $M\setminus L$. Put $u_j,v_j$ as above in $\{\tfrac12<jx<2\}\times\supp g\subset M\setminus L$. Then $u_j,v_j\in C_c^\infty(M\setminus L)\subset C^\infty(M)\subset C^{-\infty}(M,L)$ (by (7.28)), and for every $P\in\Diff(M,L)$ the function $Pu_j$ is smooth and compactly supported away from $L$, where $\xb$ is bounded below, so $\xb^{-m_1}Pu_j$ is bounded: $u_j,v_j\in J^{m_1}(M,L)$ for all $m_1$. Further $u_j,v_j\to0$ in $C^\infty(M\setminus L)$, $\sup\xb^{-m'}|u_j|\ge1$, and the same 0-neighborhood argument applies to $J(M,L)=\bigcup_kJ^{-k}(M,L)$. \textbf{Both assertions of Cor.~7.13 are false.} (This does not use the identification (7.26).)

\section*{7. Witness (f): $S^\infty(U\times\R^l)$ is not regular for non-compact $U$; Cor.~3.6, p.~18, [M] \S2.1.8}
$U\subset\R^n$ open, nonempty, $n\ge1$ (such a $U$ is never compact). Choose $x_j\in U$ with no accumulation point in $U$, and closed balls $\bar B(x_j,r_j)\subset U$, pairwise disjoint, such that every compact $K\subset U$ meets only finitely many of them (take a compact exhaustion $K_1\subset K_2\subset\cdots$ of $U$; pass to a subsequence with $x_j\notin K_j$; choose $r_j<\tfrac12\min\{\dist(x_j,K_j),\dist(x_j,\R^n\setminus U),\dist(x_j,\bar B(x_i,r_i))\ (i<j)\}$; then $\bar B(x_j,r_j)\cap K_j=\emptyset$, and $K\subset K_i$ meets no ball with $j\ge i$). $\chi_j\in C_c^\infty(B(x_j,r_j))$, $\chi_j(x_j)=1$.

$b_j(x,\xi):=\chi_j(x)(1+|\xi|^2)^{j/2}$.
\begin{enumerate}
\item $b_j\in S^j$: $|\partial^\beta_\xi(1+|\xi|^2)^s|\le C_{\beta,s}(1+|\xi|)^{2s-|\beta|}$ (induction: each derivative produces factors $\xi_i(1+|\xi|^2)^{-1}$ or $(1+|\xi|^2)^{-1}$).
\item $b_j\notin S^m$ for $m<j$: for $K\ni x_j$, $\|b_j\|_{K,0,0,m}\ge\sup_\xi(1+|\xi|^2)^{j/2}(1+|\xi|)^{-m}\ge2^{-j/2}\sup_\xi(1+|\xi|)^{j-m}=\infty$ (using $(1+|\xi|)^2\le2(1+|\xi|^2)$).
\item $\{b_j\}$ is bounded in $S^\infty$. Let $W$ be an absolutely convex 0-neighborhood of $S^\infty$ (these form a basis). For each $k\in\N$, $W\cap S^k$ is a 0-neighborhood of $S^k$, hence contains a basic set $W_k:=\{a\in S^k:\|a\|_{K_k,\alpha,\beta,k}<\varepsilon_k\text{ for }(\alpha,\beta)\in F_k\}$, $K_k\subset U$ compact, $F_k$ finite, $\varepsilon_k>0$. Choose $j_0$ with $\bar B(x_j,r_j)\cap K_1=\emptyset$ for $j\ge j_0$. Fix $\rho\in C_c^\infty(\R^l)$ with $\rho=1$ on $B(0,1)$, $\supp\rho\subset B(0,2)$, $\rho_R(\xi):=\rho(\xi/R)$, $R\ge1$. For $j\ge j_0$ split $b_j=b_j\rho_R+b_j(1-\rho_R)$.
\begin{itemize}
\item $b_j\rho_R$ is smooth with compact $\xi$-support, so it lies in $S^{-\infty}\subset S^1$, and it vanishes for $x\in K_1$. So every seminorm in $W_1$ of $2b_j\rho_R$ is $0<\varepsilon_1$: $2b_j\rho_R\in W_1\subset W$.
\item $b_j(1-\rho_R)\in S^{j+1}$ vanishes for $|\xi|\le R$. For $(\alpha,\beta)\in F_{j+1}$: $\partial^\alpha_x\partial^\beta_\xi[b_j(1-\rho_R)]=\partial^\alpha\chi_j(x)\sum_{\beta'+\beta''=\beta}C_{\beta'\beta''}\partial^{\beta'}(1+|\xi|^2)^{j/2}\partial^{\beta''}(1-\rho_R)$. The term $\beta''=0$ is supported in $|\xi|\ge R$ and bounded by $C(1+|\xi|)^{j-|\beta|}=C(1+|\xi|)^{j+1-|\beta|}(1+|\xi|)^{-1}\le C(1+|\xi|)^{j+1-|\beta|}/(1+R)$. A term with $\beta''\ne0$ is supported in $R\le|\xi|\le2R$, where $|\partial^{\beta''}(1-\rho_R)|\le C_{\beta''}R^{-|\beta''|}$, so it is bounded by $C(1+|\xi|)^{j-|\beta'|}R^{-|\beta''|}\le C(1+|\xi|)^{j-|\beta|}((1+|\xi|)/R)^{|\beta''|}\le3^{|\beta''|}C(1+|\xi|)^{j-|\beta|}\le3^{|\beta''|}C(1+|\xi|)^{j+1-|\beta|}/(1+R)$ on that region ($1+|\xi|\ge1+R$ there). Hence $\|b_j(1-\rho_R)\|_{K,\alpha,\beta,j+1}\le C_{j,\alpha,\beta}/(1+R)$ for every compact $K$. Taking $R=R_j$ large, $2b_j(1-\rho_R)\in W_{j+1}\subset W$.
\item Absolute convexity: $b_j=\tfrac12(2b_j\rho_R)+\tfrac12(2b_j(1-\rho_R))\in W$.
\end{itemize}
For the finitely many $j<j_0$, $b_j\in\lambda_jW$ since $W$ is absorbing. So $\{b_j\}\subset\lambda W$ with $\lambda:=\max\{1,\lambda_j\}$: \textbf{bounded in $S^\infty$.}
\item By 2, $\{b_j\}$ is contained in no step $S^m$. So $(S^m)_m$ is \textbf{not regular}; by the definition on p.~4 it is not boundedly retractive, and by ``[39, Corollary 6.5]'' as quoted on p.~5 (acyclic $\Rightarrow$ regular) it is \textbf{not acyclic}. \textbf{Cor.~3.6's acyclicity and bounded-retractivity clauses are false for every non-compact $U$}, i.e.\ for every nonempty open $U\subset\R^n$, $n\ge1$. The printed derivations of the completeness and Montel clauses go through these properties, so they are unproved as printed (the note does not claim they are false). For $U=\R^0=\{0\}$, i.e.\ $S^\infty(\R^l)$, Cor.~3.6 is true by \S9 below ($d=0$, $K=\{0\}$).
\item Direct failure of the criterion (p.~4) for the same spectrum, for completeness of the record: a 0-neighborhood $V$ of $S^m$ controls finitely many seminorms, hence one compact $K$; for $x_0\in U\setminus K$ take $\psi\in C_c^\infty(U\setminus K)$ with $\psi(x_0)=1$, $\varphi\in C_c^\infty(\{\tfrac12<|\eta|<2\})$ with $\varphi(e_1)=1$, and the dilated annular bumps $a_j:=\psi(x)\,j^{m'}\varphi(\xi/j)$. Then $a_j\in S^{-\infty}\subset S^m$, every $K$-seminorm of $a_j$ vanishes, so $a_j\in V$ for all $j$; $\partial^\alpha_x\partial^\beta_\xi a_j=\partial^\alpha\psi(x)\,j^{m'-|\beta|}(\partial^\beta\varphi)(\xi/j)$ is supported in $\{j/2<|\xi|<2j\}$, where $1+|\xi|\asymp j$, so every (3.1)-seminorm of $a_j$ of order $m''$ is $\le C_{\alpha,\beta}\,j^{m'-|\beta|}j^{|\beta|-m''}=C_{\alpha,\beta}\,j^{m'-m''}\to0$: $a_j\to0$ in $S^{m''}$; and for $K'\ni x_0$, $\|a_j\|_{K',0,0,m'}\ge j^{m'}\varphi(e_1)(1+j)^{-m'}=(j/(1+j))^{m'}\to1$ (at $x=x_0$, $\xi=je_1$): $a_j\not\to0$ in $S^{m'}$. So the $m'$- and $m''$-topologies differ on every 0-neighborhood of $S^m$, for every $m<m'<m''$.
\item The extension on p.~18 (``Propositions 3.2 and 3.3 and Corollaries 3.4 to 3.6 can be directly extended to this setting'', symbols on a vector bundle $E$ over a manifold $M$): Prop.~3.2 and Cor.~3.4 fail there for every $M$ (witness (a) in a chart); Cor.~3.6 fails for non-compact $M$ (witness (f) in a chart, with a locally finite family of balls in $M$). [M] \S2.1.8 (arXiv p.~15) restates Cors.~3.4--3.6 for arbitrary open $U\subset\R^n$: the coincidence sentence is false for every $U$, the acyclic/boundedly retractive sentence for every non-compact $U$ ($n\ge1$).
\end{enumerate}

\section*{8. Consumers of the false statements inside [P] (read on the cited pages)}
\begin{itemize}
\item Cor.~3.5 (p.~17): proof cites Cor.~3.4; statement true (\S10.3 below).
\item Cor.~3.6 (p.~18), Cor.~4.7 (p.~23), Cor.~6.16 (p.~39), Cor.~6.22 (p.~40), Cor.~6.28 (p.~41), Cor.~7.15 (p.~57), Cor.~7.23 (p.~58): the acyclicity clauses are derived from Cor.~3.4 / 4.5 / 6.14 / 6.21 / 6.27 / 7.13 / 7.22; true for compact bases with the proofs of \S9--\S10.
\item Remark~3.8 (p.~18): uses sequential retractivity of $S^\infty(U\times\R^l)$ from Cor.~3.6; valid for compact base ($S^\infty(\R^l)$); for non-compact $U$ the remark's argument needs the compactly based version.
\item Claim~6.46 (pp.~48--49): ``By Corollary 6.14, there is some 0-neighborhood $V\subset A(M)$ such that $V\cap A^m(M)=W\cap A^m(M)$''; repaired in \S10.1.
\item Prop.~8.8 (p.~64): ``formally the same proof, using Proposition 7.29 and Corollaries 7.13, 7.15 and 7.31''; repaired in \S10.2.
\item Cor.~6.27 (p.~41) and Cor.~7.22 (p.~58): proofs ``formally the same'' as Cor.~6.21 (which rests on Cor.~4.5); statements true (\S10.4).
\item Cor.~6.21 (p.~40) is itself ``a consequence of Corollary 4.5 applied to $(\Mb,\partial M)$''; false (\S5).
\end{itemize}

\section*{9. What survives: the compactly based models and the acyclicity criterion}
\subsection*{9.1 The models}
For $K\subset\R^d$ compact ($d\ge0$), $S^m_K:=\{a\in C^\infty(\R^d\times\R^l):a(x,\xi)=0\text{ for }x\notin K,\ N_m(a;\gamma)<\infty\ \forall\gamma\}$, $N_m(a;\gamma):=\sup_{x,\xi}|\partial^\gamma a(x,\xi)|(1+|\xi|)^{|\beta|-m}$, $\gamma=(\alpha,\beta)$. Fr\'echet (weighted $C^\infty$ space; Cauchy sequences converge locally uniformly with all derivatives and the bounds pass to the limit). If $K\subset U$ open, $S^m_K$ is the closed subspace of $S^m(U\times\R^l)$ of symbols vanishing for $x\notin K$ (extension by zero is smooth, since such a symbol vanishes on $(U\setminus K)\times\R^l$), with the induced topology (a seminorm over compact $K'\subset U$ equals the seminorm over $K'\cap K\le N_m$; $N_m$ is the seminorm over $K$). The (4.10)-symbol $a_j$ of $f_ju$ vanishes for $x''$ outside the compact projection $K_j$ of $\supp f_j$, so (4.10) maps $I^m(M,L)$ into $C^\infty(M\setminus L)\oplus\bigoplus_jS^{\bar m}_{K_j}$, with the same map for every $m$.

\subsection*{9.2 Lemma A (one variable; Vandermonde)}
$N\ge1$, $f\in C^N([t,t+Nh])$, $h>0$. Then for $0\le j\le N$: $|f^{(j)}(t)|\le C_N(h^{-j}\|f\|_\infty+h^{N-j}\|f^{(N)}\|_\infty)$ (sup norms on $[t,t+Nh]$).

\noindent Proof. Taylor at $t$: $f(t+kh)=\sum_{i=0}^{N-1}y_ik^i+r_k$ with $y_i:=f^{(i)}(t)h^i/i!$ and $|r_k|\le(kh)^N\|f^{(N)}\|/N!\le(Nh)^N\|f^{(N)}\|/N!$, $k=0,\dots,N-1$. The matrix $(k^i)_{0\le k,i\le N-1}$ is Vandermonde with distinct nodes $0,\dots,N-1$, hence invertible with an inverse depending only on $N$. So $|y_j|\le C_N\max_k|f(t+kh)-r_k|\le C_N(\|f\|+(Nh)^N\|f^{(N)}\|/N!)$, and $|f^{(j)}(t)|=j!h^{-j}|y_j|\le C'_N(h^{-j}\|f\|+h^{N-j}\|f^{(N)}\|)$. (For $j=N$ use $f^{(N)}$ directly.)\qed

\subsection*{9.3 Lemma B (box, per-variable scales)}
$F\in C^\infty$ on a neighborhood of the box $Q=\prod_{i=1}^D[t_i,t_i+Nh_i]$, $\gamma\in\N^D$ with $\gamma_i\le N$. Then
\[
|\partial^\gamma F(t)|\le C_{N,D}\sum_{S\subset\{1,\dots,D\}}\Big(\prod_{i\in S}h_i^{N-\gamma_i}\Big)\Big(\prod_{i\notin S}h_i^{-\gamma_i}\Big)\sup_Q|\partial^{N\chi_S}F|,
\]
where $N\chi_S$ is the multi-index with entry $N$ at $i\in S$ and $0$ elsewhere.

\noindent Proof. Induction on $D$. $D=1$ is Lemma A. For $D>1$ write $\gamma=(\gamma_1,\tilde\gamma)$. Lemma A in the first variable applied to $s\mapsto\partial^{\tilde\gamma}F(s,\tilde t)$: $|\partial^\gamma F(t)|\le C(h_1^{-\gamma_1}\sup_s|\partial^{\tilde\gamma}F(s,\tilde t)|+h_1^{N-\gamma_1}\sup_s|\partial_1^N\partial^{\tilde\gamma}F(s,\tilde t)|)$, $s\in[t_1,t_1+Nh_1]$. For each $s$ apply the $(D-1)$-case to $F(s,\cdot)$ and to $\partial_1^NF(s,\cdot)$ at $\tilde t$ on the box $\tilde Q$: the first gives the terms with $1\notin S$, the second the terms with $1\in S$.\qed

\subsection*{9.4 Proposition (interpolation inequality)}
Let $m<m'<m''$ and $\gamma=(\alpha,\beta)$ with $|\gamma|\ge1$. Put $c:=(m'-m)/(2|\gamma|)$, choose $N\in\N$ with $N>|\gamma|+(m''-m')/c$, $R_0:=\max\{1,(2N\sqrt l)^{1/c}\}$, $\Gamma:=\{\gamma\}\cup\{N\chi_S:\emptyset\ne S\subset\{1,\dots,d+l\}\}$. There is $C$ (depending on $m,m',m'',\gamma,d,l$, not on $K$) such that for all $a\in S^{m''}_K$ and $R\ge R_0$:
\begin{equation}\label{interp}
N_{m'}(a;\gamma)\ \le\ C\Big[R^{-(m'-m)/2}\,N_m(a;0)+R^{\,m''-m'}\max_{\gamma'\in\Gamma}N_{m''}(a;\gamma')\Big].
\end{equation}
For $\gamma=0$: $N_{m'}(a;0)\le R^{m''-m'}N_{m''}(a;0)+R^{m-m'}N_m(a;0)$, which is \eqref{interp} with $\Gamma=\{0\}$, $C=1$, $R_0=1$, since $m-m'<-(m'-m)/2$.

\noindent Proof. Fix $(x,\xi)$, $w:=1+|\xi|$.

(i) $w\le R$: $|\partial^\gamma a|w^{|\beta|-m'}=|\partial^\gamma a|w^{|\beta|-m''}\cdot w^{m''-m'}\le R^{m''-m'}N_{m''}(a;\gamma)$. (For $\gamma=0$ the same.)

(ii) $w>R$: box $Q:=\prod_{i\le d}[x_i,x_i+Nh_x]\times\prod_{i\le l}[\xi_i,\xi_i+Nh_\xi]$, $h_x:=w^{-c}$, $h_\xi:=w^{1-c}$. On $Q$, $|\xi'-\xi|\le\sqrt l\,Nw^{1-c}\le w/2$ because $w\ge R_0\ge(2N\sqrt l)^{1/c}$; so $1+|\xi'|\in[w/2,3w/2]$. Hence $\sup_Q|a|\le N_m(a;0)\sup_Q(1+|\xi'|)^m\le2^{|m|}w^mN_m(a;0)$, and $\sup_Q|\partial^{N\chi_S}a|\le N_{m''}(a;N\chi_S)\sup_Q(1+|\xi'|)^{m''-N|S_\xi|}\le2^{|m''|+Nl}w^{m''-N|S_\xi|}N_{m''}(a;N\chi_S)$, where $S_\xi:=S\cap\{\xi\text{-indices}\}$, $S_x:=S\cap\{x\text{-indices}\}$. Lemma B ($D=d+l$, $\gamma_i\le|\gamma|<N$) and multiplication by $w^{|\beta|-m'}$:
\begin{itemize}
\item $S=\emptyset$: $\prod h_i^{-\gamma_i}=w^{c|\alpha|}w^{-(1-c)|\beta|}$; total exponent $c|\alpha|-(1-c)|\beta|+m+|\beta|-m'=c|\gamma|+m-m'=-(m'-m)/2$. Contribution $\le Cw^{-(m'-m)/2}N_m(a;0)\le CR^{-(m'-m)/2}N_m(a;0)$.
\item $S\ne\emptyset$: $(\prod_{i\in S}h_i^{N-\gamma_i})(\prod_{i\notin S}h_i^{-\gamma_i})=w^{-c(N|S_x|-|\alpha_S|)}w^{(1-c)(N|S_\xi|-|\beta_S|)}w^{c|\alpha_{S^c}|}w^{-(1-c)|\beta_{S^c}|}$; with the sup factor $w^{m''-N|S_\xi|}$ and the weight $w^{|\beta|-m'}$ the total exponent is $-cN|S_x|+c|\alpha|+N|S_\xi|-cN|S_\xi|-|\beta|+c|\beta|+m''-N|S_\xi|+|\beta|-m'=-c(N|S|-|\gamma|)+m''-m'\le-c(N-|\gamma|)+m''-m'<0$ by the choice of $N$. Contribution $\le CN_{m''}(a;N\chi_S)\le CR^{m''-m'}N_{m''}(a;N\chi_S)$ ($R\ge1$).
\end{itemize}
Taking the sup over $(x,\xi)$ gives the inequality. For $\gamma=0$ and $w>R$: $|a|w^{-m'}=|a|w^{-m}w^{m-m'}\le R^{m-m'}N_m(a;0)$.\qed

\noindent Where compact base support enters: $N_m(a;0)$ and $N_{m''}(a;\cdot)$ are global suprema, so the box may be placed anywhere; for $S^m(U\times\R^l)$ with non-compact $U$ the seminorms over a fixed compact $K$ control nothing outside $K$, which is exactly witness (f), item 5.

\subsection*{9.5 From the inequality to the criterion of p.~4}
\Mstar\ at $(m,m',m'')$: $V:=\{a\in S^m_K:N_m(a;0)<1\}$. Given a basic 0-neighborhood $O=\{N_{m'}(\cdot;\gamma_i)<\varepsilon,\ i\le p\}$ of $S^{m'}_K$, let $C_i$, $R_{0,i}$, $\Gamma_i$ be the constants of 9.4 for $\gamma_i$; choose $R\ge\max_iR_{0,i}$ with $C_iR^{-(m'-m)/2}<\varepsilon/2$ for all $i$, then $\delta>0$ with $C_iR^{m''-m'}\delta<\varepsilon/2$ for all $i$, and $W:=\{N_{m''}(\cdot;\gamma')<\delta,\ \gamma'\in\bigcup_i\Gamma_i\}$. By 9.4, $V\cap W\subset O$.

\Mstar\ $\Rightarrow$ coincidence on $V$ (the form quoted on p.~4): the $S^{m''}$-topology is coarser than the $S^{m'}$-topology, so it suffices that for $a\in V$ and $O$ there is $W_0$ with $(a+W_0)\cap V\subset a+O$. Take $W_0$ with $V\cap W_0\subset\tfrac12O$ ($\tfrac12O$ is also a 0-neighborhood). If $b\in V$, $b-a\in W_0$, then $(b-a)/2\in V$ (absolutely convex) and $(b-a)/2\in W_0$, so $(b-a)/2\in\tfrac12O$, $b-a\in O$.\qed\ So for every $m$ the topologies of $S^{m'}_K$ and $S^{m''}_K$ coincide on the 0-neighborhood $V$ of $S^m_K$, for all $m'>m$ and $m''>m'$: this is ``[39, Theorem 6.1]'' as quoted, with $k'=k+1$; hence \textbf{$\bigcup_mS^m_K$ is acyclic} (indices in $\N$ by cofinality, p.~5).

Transfer rules (each immediate): \Mstar\ holds for a constant spectrum $X_k=X$ ($V=X$, $W=O$); for finite products (product neighborhoods); for a spectrum $(Y_k)$ with $Y_k$ carrying the initial topology of a fixed linear map $j$ with $j(Y_k)\subset X_k$ ($V_Y:=j^{-1}(V)\cap Y_k$; for $O_Y\supseteq j^{-1}(O)\cap Y_{k'}$ take $W_Y:=j^{-1}(W)\cap Y_{k''}$; then $V_Y\cap W_Y\subset Y_k\subset Y_{k'}$ and $j(V_Y\cap W_Y)\subset V\cap W\subset O$); in particular for closed subspaces with the induced topologies.

\subsection*{9.6 The $b$-collar model}
$B^m_K:=\{v\in C^\infty([\varrho_0,\infty)\times\R^d):v=0\text{ for }y\notin K,\ N^b_m(v;\gamma):=\sup e^{m\varrho}|\partial^\gamma v|<\infty\ \forall\gamma\}$ with $\varrho=-\log x$ ($x=e^{-\varrho}$, $x\partial_x=-\partial_\varrho$, $x^{-m}=e^{m\varrho}$); $B^m\subset B^{m'}$ for $m'<m$, matching (6.38). Write the three orders as $-k>-k'>-k''$ with $k<k'<k''$, so that the steps are $B^{-k}_K\subset B^{-k'}_K\subset B^{-k''}_K$ with weights $w^{-k}$, $w^{-k'}$, $w^{-k''}$, $w:=e^\varrho$. The proof of 9.4 applies with $w$ in place of $1+|\xi|$ and with $(k,k',k'')$ in the roles of $(m,m',m'')$, i.e.\ with $c:=(k'-k)/(2|\gamma|)>0$ and $N>|\gamma|+(k''-k')/c$ (not with $m=-k$ etc.\ substituted literally, which would give $c<0$); all box sides are $h:=w^{-c}$ (the $\varrho$-box $[\varrho,\varrho+Nh]\subset[\varrho_0,\infty)$ and, for $w\ge N^{1/c}$, $e^{\varrho'}\in[w,ew]$ on $Q$), and there is no gain from $\varrho$-derivatives: $\sup_Q|v|\le e^kw^kN^b_{-k}(v;0)$, $\sup_Q|\partial^{N\chi_S}v|\le e^{k''}w^{k''}N^b_{-k''}(v;N\chi_S)$; the $S=\emptyset$ exponent is $c|\gamma|+k-k'=-(k'-k)/2$, the $S\ne\emptyset$ exponent is $-c(N|S|-|\gamma|)+k''-k'<0$, and the regime $w\le R$ gives $R^{k''-k'}N^b_{-k''}(v;\gamma)$. Hence $N^b_{-k'}(v;\gamma)\le C[R^{-(k'-k)/2}N^b_{-k}(v;0)+R^{k''-k'}\max_{\gamma'\in\Gamma}N^b_{-k''}(v;\gamma')]$ for $R\ge R_0$, and the same conclusion: \Mstar\ for $(B^{-k}_K)_k$, acyclic.

\subsection*{9.7 Assembly ($M$ compact, or compact with boundary)}
\begin{itemize}
\item $I(M,L)$: $I^m$ carries the initial topology of (4.10) into $C^\infty(M\setminus L)\times\prod_jS^{\bar m}_{K_j}$ (9.1), the same map for all $m$; the target spectrum is a finite product of a constant spectrum and models satisfying \Mstar\ at every triple; so $(I^m)_m$ satisfies \Mstar, and $\bigcup_mI^m$ is acyclic. By (4.12) the spectra $(I^m)_m$ and $(I^{(s)})_s$ are mutually cofinal with continuous inclusions, so the LF-space $I(M,L)$ is the same and acyclicity is its property (``[39, Chapter 6, p.~111]'', p.~5). Steps are Fr\'echet (p.~22).
\item $A(M)$: take a partition of unity $\{\lambda_j\}\cup\{\mu\}$ with $\lambda_j$ supported in collar charts and $\mu$ in $\Mo$, and $\Lambda\colon u\mapsto((\lambda_ju)(e^{-\varrho},y))_j\oplus\mu u\in\prod_jB^m_{K_j}\times C^\infty_{K'}(\Mo)$ ($K_j$ the $y$-projection of $\supp\lambda_j$, $K'=\supp\mu$). Continuity: $\sup e^{m\varrho}|\partial^a_\varrho\partial^b_y(\lambda_ju)|=\sup x^{-m}|(x\partial_x)^a\partial_y^b(\lambda_ju)|$ is the seminorm of $P:=(x\partial_x)^a\partial_y^b\circ\lambda_j\in\Diff_b(M)$; $\mu u$ by (6.33). Initial topology: for $P\in\Diff_b(M)$, $Pu=\sum_jP(\lambda_ju)+P(\mu u)$; on the closure of a collar chart $P$ is a $C^\infty$-combination of the monomials $(x\partial_x)^a\partial_y^b$ ($b$-vector fields are $C^\infty$-combinations of $x\partial_x$, $\partial_y$) with bounded coefficients, so $\sup x^{-m}|P(\lambda_ju)|\le C\sum N^b_m((\lambda_ju)(e^{-\varrho},y);(a,b))$; and $\sup x^{-m}|P(\mu u)|\le C\|\mu u\|_{C^k(K')}$ since $x$ is bounded below on $K'$. $\Lambda$ is $m$-independent. So $(A^m)_m$ satisfies \Mstar\ and $A(M)$ is acyclic. (Injectivity of $\Lambda$, not needed for \Mstar: $u=\sum\lambda_ju+\mu u$ on $\Mo$, and elements of $C^{-\infty}(M)=\dot C^\infty(M;\Omega)'$ are determined on $\Mo$ since $C_c^\infty(\Mo;\Omega)$ is dense in $\dot C^\infty(M;\Omega)$.)
\item $\dot A(M)$: $\dot A^m=I^m_M(\Mb,\partial M)$ closed subspace of $I^m(\Mb,\partial M)$ (6.47): transfer rule. $\mathcal K(M)$: $\mathcal K^m=$ null space of $R$ restricted to $\dot A^m$, closed subspace (p.~41). $J(M,L)$: $J^m\cong A^m(\MB)$ via (7.26) by definition (p.~56). $K(M,L)$: $K^m(M,L)=I^m_L(M,L)$ closed subspace of $I^m(M,L)$ (p.~58).
\item Hence all six LF-spaces are acyclic. Their steps are Fr\'echet spaces: $I^m(M,L)$ by p.~22; $A^m(M)$ by Cor.~6.38 with Remark~6.41 (\S\S6.19--6.21 of [P] cite none of the statements under repair), or directly: the topology of $A^m(M)$ is given by the countable family (6.41), so for a Cauchy sequence $(u_\nu)$ each $Pu_\nu$ is Cauchy in the Banach space $x^mL^\infty(M)$, which sits continuously in $C^{-\infty}(M)$ (p.~37); with $P=1$ this gives $u_\nu\to u$ in $C^{-\infty}(M)$, continuity of $P$ on $C^{-\infty}(M)$ gives $Pu=\lim Pu_\nu\in x^mL^\infty(M)$ for every $P$, so $u\in A^m(M)$ and $u_\nu\to u$ there; $J^m(M,L)\cong A^m(\MB)$ by (7.26) inherits it; $\dot A^m(M)$, $\mathcal K^m(M)$, $K^m(M,L)$ are closed subspaces of Fr\'echet spaces. (The models $S^m_K$ and $B^m_K$ need no such statement: the transfer rule uses only \Mstar\ of the target spectrum.) Therefore the six LF-spaces are boundedly/compactly/sequentially retractive, complete, regular (``[39, Proposition 6.4], [39, Corollary 6.5]'', p.~5). This proves the acyclicity clauses of Cor.~3.6 (compact base), 4.7, 6.16, 6.22, 6.28, 7.15, 7.23.
\item Montel clauses: bounded subsets of $S^m_K$ are relatively compact in $S^{m'}_K$ for $m'>m$ (Arzel\`a--Ascoli gives a $C^\infty$-convergent subsequence on compacta; the limit $a$ satisfies the $S^m$ bounds; $N_{m'}(a_\nu-a;\gamma)\le2C(1+R)^{m-m'}$ on $|\xi|>R$ plus a uniform-convergence term on $K\times\bar B(0,R)$). For a closed bounded $B\subset I(M,L)$: by bounded retractivity $B$ is bounded in some $I^m$ with the $I(M,L)$- and $I^{m'}$-topologies agreeing on $B$ ($m'>m$); its (4.10)-image is bounded in $C^\infty(M\setminus L)\times\prod S^{\bar m}_{K_j}$, hence relatively compact in the $\bar m'$-level; $B$ is closed in $I^{m'}$ ($I^{m'}\to I(M,L)$ continuous) and $I^{m'}$ is closed in the product (complete), so $B$ is compact in $I^{m'}$, hence in $I(M,L)$. With barreledness (Cor.~4.2), $I(M,L)$ is Montel. Same for the other five spaces through the same embeddings. (The printed Montel step of Cor.~4.7, p.~23, cites Cor.~3.6 for $S^\infty(N^*L_j;\Omega N^*L_j)$ over the non-compact base $L_j$; the compactly based models avoid this as well.)
\item The inclusions $S^m_K\to S^{m'}_K$ are \textbf{not} compact operators (take $K$ with $\interior K\ne\emptyset$, since $S^m_K=\{0\}$ otherwise, and $\theta_1\in C_c^\infty(\interior K)$, $\theta_1\not\equiv0$): a 0-neighborhood $\{N_m(\cdot;\gamma)<1,\ \gamma\in F\}$ ($F$ finite, orders $\le k_0$ in $x$) contains $a_\lambda:=\varepsilon\lambda^{-k_0}\sin(\lambda x_1)\theta_1(x)\theta(\xi)$ for all $\lambda\ge1$ ($\varepsilon$ small), while $N_{m'}(a_\lambda;(\alpha_0,0))\gtrsim\varepsilon\lambda\to\infty$ for $\alpha_0=(k_0+1)e_1$. So the p.~5 remark (``compact $\Rightarrow$ clearly acyclic'') does \textbf{not} apply to these spectra; the interpolation route is needed.
\end{itemize}

\section*{10. Repairs of the two non-acyclicity consumers, density, and the two true $K$-statements}
\subsection*{10.1 Claim~6.46 (p.~48)}
Given bounded $A\subset A(M)$; replace $A$ by its absolutely convex hull (still bounded, contains $0$). Bounded retractivity (9.7): $A\subset A^m(M)$ bounded and the $A(M)$-topology equals the $A^m$-topology on $A$; then for $m'<m$ ($A^m\subset A^{m'}$, (6.38)) the three topologies $A(M)\le A^{m'}\le A^m$ coincide on $A$. Let $E_{m'}\colon A^{m'}(M)\to\dot A^{(s)}(M)$ be the partial extension map of Prop.~6.29 ($RE_{m'}=1$ on $A^{m'}$). $B:=E_{m'}(A)$ is bounded in $\dot A^{(s)}$ hence in $\dot A(M)$. Given a 0-neighborhood $U$ of $\dot A(M)$, $W:=E_{m'}^{-1}(U\cap\dot A^{(s)})$ is a 0-neighborhood of $A^{m'}$; since the topologies coincide on $A$, there is a 0-neighborhood $V$ of $A(M)$ with $V\cap A\subset W\cap A$; for $a\in A\cap V$, $E_{m'}a\in B\cap U$ and $a=R(E_{m'}a)\in R(B\cap U)$.\qed\ Differences from the printed proof: ``$V\cap A^m=W\cap A^m$'' (false in general, \S6) is replaced by ``$V\cap A\subset W$''; and the index must satisfy $m'<m$, since $E_{m'}$ is defined on $A^{m'}\supset A^m$ (the printed text has $m'>m$ with $E_{m'}$ on $A^m$, which does not match Prop.~6.29's domain).

\subsection*{10.2 Prop.~8.8 (p.~64)}
Same lemma with $X_m:=J^m(M,L)$, $Y:=I(M,L)$, $Y_s:=I^{(s)}(M,L)$, $R=(7.31)$, $E_{m'}$ from Cor.~7.31 with $m'<m$ (as in 10.1: $E_{m'}$ is defined on $J^{m'}(M,L)\supset J^m(M,L)$, so the index must satisfy $m'<m$, not $m'>m$; $RE_{m'}=1$ on $J^{m'}$), bounded retractivity of $J(M,L)$ from 9.7. The printed ``using \dots\ Corollaries 7.13, 7.15 and 7.31'' needs only Cor.~7.15 (acyclicity, repaired) and Cor.~7.31.

\subsection*{10.3 Cor.~3.5 (p.~17), directly}
$a\in S^m$, $m<m'$, $\chi\in C_c^\infty(\R^l)$, $\chi=1$ on $B(0,1)$; $a_R:=\chi(\xi/R)a$. Then $a-a_R=(1-\chi(\xi/R))a$ vanishes for $|\xi|\le R$, and by Leibniz, on $|\xi|\ge R$, $|\partial^\alpha_x\partial^\beta_\xi[(1-\chi(\xi/R))a]|\le C(1+|\xi|)^{m-|\beta|}$ (a derivative of order $\beta'\ne0$ of $1-\chi(\xi/R)$ is supported where $1+|\xi|\asymp R$ and is $O(R^{-|\beta'|})=O((1+|\xi|)^{-|\beta'|})$); hence $\|a-a_R\|_{K,\alpha,\beta,m'}\le C\sup_{|\xi|\ge R}(1+|\xi|)^{m-m'}=CR^{m-m'}\to0$. Then $\phi a_R\in C_c^\infty$ for $\phi\in C_c^\infty(U)$, $\phi=1$ near $K$, and all $K$-seminorms of $a_R-\phi a_R$ vanish. So $C_c^\infty$ is dense in $S^m$ for the $S^{m'}$-topology; the second assertion follows as printed. Cor.~4.6's printed proof uses only Cor.~3.5 and Prop.~4.3 (p.~23). (Remark~3.7, p.~18, anticipates a direct proof of Cor.~3.5.)

\subsection*{10.4 Cor.~6.27 (p.~41) and Cor.~7.22 (p.~58) are true}
By Prop.~7.26 (p.~59), $K(M,L)\cong\bigoplus_{k\ge0}C^\infty(L;\Omega^{-1}NL)$ (locally convex direct sum), $u\mapsto\sum_k\partial_x^k\delta_L\otimes v_k$, a TVS-isomorphism proved in [P] from Prop.~7.7 and the open mapping theorem, not from any coincidence statement. So every $u\in K(M,L)$ is a finite sum of Dirac layers. In an adapted chart $(x,y)$, $L=\{x=0\}$, the partial Fourier transform (4.8) of $\partial_x^k\delta(x)v(y)$ is $(i\xi)^kv(y)$: a polynomial in $\xi$ of degree $k$. For $f\in C^\infty$, $f\cdot\partial_x^k\delta=\sum_{i\le k}\binom ki(-1)^{k-i}(\partial_x^{k-i}f)|_{x=0}\,\partial_x^i\delta$ (pair with a test function and expand $(f\phi)^{(k)}(0)$), so the symbols of $f_ju$ are again polynomials of degree $\le k_0:=$ max layer index, with $C_c^\infty$ coefficients. A polynomial symbol $\sum_{k\le k_0}(i\xi)^kw_k(y)$ lies in $S^{\bar m}$ iff $w_k=0$ for $k>\bar m$ (if $w_{k_1}\not\equiv0$ with $k_1>\bar m$ maximal, $|a|(1+|\xi|)^{-\bar m}\sim|w_{k_1}||\xi|^{k_1-\bar m}\to\infty$). Hence $K^m(M,L)=\{$finite sums with $k\le\bar m\}$. On $P_{k_0}:=\{\sum_{k\le k_0}(i\xi)^kw_k:w_k\in C_c^\infty(K)\}$ and $\bar m_1\ge k_0$: $N_{\bar m_1}(a;\alpha,\beta)\le\sum_k\|\partial^\alpha w_k\|_\infty\sup_\xi|\partial_\xi^\beta\xi^k|(1+|\xi|)^{|\beta|-\bar m_1}\le C\sum_k\|\partial^\alpha w_k\|_\infty$ ($k-|\beta|+|\beta|-\bar m_1\le0$); conversely $\partial_\xi^{k_0}a=i^{k_0}k_0!\,w_{k_0}(y)$, so $\|\partial^\alpha w_{k_0}\|_\infty\le C|\partial^\alpha\partial_\xi^{k_0}a|_{\xi=0}\le CN_{\bar m_1}(a;\alpha,k_0)$, and peeling off the top term (whose seminorms are $\le C\|\partial^\alpha w_{k_0}\|$) one obtains $\|\partial^\alpha w_k\|_\infty\le C\max_{\beta\le k_0}N_{\bar m_1}(a;\alpha,\beta)$ for all $k$. So on $P_{k_0}$ every $S^{\bar m_1}$-topology with $\bar m_1\ge k_0$ equals the $C^\infty$-topology of the coefficient tuple. For $m<m',m''$: $k_0\le\lfloor\bar m\rfloor\le\bar m<\bar m',\bar m''$, so the $I^{m'}$- and $I^{m''}$-topologies restricted to $K^m(M,L)$ (initial topologies of (4.10), whose $C^\infty(M\setminus L)$-component is $0$ on $K$) are both the coefficient topology: \textbf{Cor.~7.22 holds.} Cor.~6.27: by (6.47) and the definitions on pp.~40 and 58, $\mathcal K^m(M)=\dot A^m_{\partial M}(M)=I^m_{\partial M}(\Mb,\partial M)=K^m(\Mb,\partial M)$ with the same subspace topologies, so Cor.~6.27 is Cor.~7.22 for $(\Mb,\partial M)$. ((6.49) says the same, but it is printed as a consequence of Cor.~6.20, hence of Cor.~6.16, so it is not used here.) \textbf{Cor.~6.27 holds.} Their printed proofs (``formally the same proofs, using Corollaries 6.21, 6.22 and 6.26'', p.~41; ``consequences of \dots\ Corollaries 6.26 to 6.28'', p.~58) go through Cor.~6.21, i.e.\ Cor.~4.5.

\section*{11. The memoir [M]}
\S2.1.8 (arXiv p.~15): the coincidence sentence is false for every $U$ (\S4), the acyclicity/retractivity sentence for every non-compact $U$ (\S7). \S2.2.2 (arXiv p.~21) cites [P, Cors.~4.2, 4.7] for $I(M,L)$ acyclic, $M$ compact: true by 9.7. \S2.5.10 (arXiv p.~38) and \S2.6.7 (arXiv p.~53) restate the second assertions of Cors.~6.14 and 7.13: false (\S6). \S5.2.1 (arXiv p.~119), \S\S5.5.3--5.5.4 (arXiv p.~122) use compact retractivity of $I(\mathcal F)=I(M,M^0;\Lambda\mathcal F)$ and $J(\mathcal F)$ with $M$ closed: these are the paper's $I(M,L)$ and $J(M,L)$ with bundle coefficients (finitely many copies of the scalar case via a local trivialization), covered by 9.7. The memoir is now in print as \cite{ALKL24m} (Lecture Notes in Mathematics 2387); the sections cited belong to its Chapters~2 and 5, and the restatements of \S\S2.1.8, 2.5.10 and 2.6.7 are presumably in print there. Beyond these pages the memoir was not examined.

\section*{12. Public record}
As of September 6, 2026 the author found no erratum: arXiv lists versions v1 to v3 only, the Crossref record of the DOI carries no correction and the publisher's Crossmark record reads ``Document is current'', and zbMATH 7901419 (Zbl 1564.46031) carries no review or corrigendum. The DOI-based citation indexes (Crossref, OpenAlex, Semantic Scholar) list no citing work; Google Scholar lists three, among them Chapter~2 of \cite{ALKL24m}, none of them a correction; the restatements of the affected sentences noted in \S11 are presumably now in print in that chapter.

\vspace{-0.5em}
\begin{thebibliography}{W}\setlength{\itemsep}{0pt}\footnotesize
\bibitem[P]{ALKL24} J.~A. \'Alvarez L\'opez, Y.~A. Kordyukov, E. Leichtnam, Topology of the space of conormal distributions, \emph{J. Pseudo-Differ. Oper. Appl.} 15 (2024), article 47, 68 pp., \url{https://doi.org/10.1007/s11868-024-00617-y}; arXiv:2304.00798 (v3, 1 June 2024).
\bibitem[M]{ALKL24m} J.~A. \'Alvarez L\'opez, Y.~A. Kordyukov, E. Leichtnam, \emph{A Trace Formula for Foliated Flows}, Lecture Notes in Mathematics, vol.~2387, Springer, Cham (2026), \url{https://doi.org/10.1007/978-3-032-15413-2}; arXiv:2402.06671 (v2, 13 February 2024). Sections are cited by number; the page numbers given are those of the arXiv version, the same in v1 and v2.
\bibitem[W]{Wen03} J. Wengenroth, \emph{Derived Functors in Functional Analysis}, Lecture Notes in Mathematics, vol.~1810, Springer, Berlin (2003).
\end{thebibliography}
\end{document}
